\chapter{KẾT LUẬN VÀ HƯỚNG PHÁT TRIỂN}
\label{chuong6}
\label{ch:ket_luan}

\vspace{-1.5cm}

\begin{quotation}
    \noindent
    \small\itshape
    Chương cuối tổng kết những đóng góp chính của khóa luận thông qua việc trả lời lần lượt ba câu hỏi nghiên cứu được đặt ra ở Chương~\ref{ch:tong_quan}, phân tích bốn nhóm hạn chế còn tồn tại, và đề xuất sáu hướng phát triển khả thi trong tương lai. Phần cuối cùng là lời kết về vai trò của các mô hình học sâu được tối ưu hóa bằng NAS trong việc giải quyết bài toán phát hiện ranh giới cảnh quay trên các nền tảng video ngắn hiện đại.
\end{quotation}
\vspace{1em}

\section{Kết luận chung}
\label{sec:ket_luan}

Khóa luận đã hoàn thành các mục tiêu nghiên cứu đề ra, tập trung vào bài toán phát hiện ranh giới cảnh quay (Shot Boundary Detection -- SBD) trong bối cảnh đặc thù của video ngắn. Phần này tổng kết kết quả nghiên cứu thông qua hai góc nhìn: trả lời trực tiếp ba câu hỏi nghiên cứu RQ1--RQ3 đã đặt ra ở Mục~\ref{sec:research_questions}, sau đó tổng hợp ba nhóm đóng góp chính.

\subsection{Trả lời các câu hỏi nghiên cứu}

\textbf{RQ1 -- NAS có giúp cải thiện hiệu năng SBD trên video ngắn so với mạng xương sống cố định?}
Câu trả lời là \textbf{có}. Trên tập SHOT, AutoShot đạt F1 = \AutoShotOriginalShotFOne, vượt mô hình cơ sở TransNetV2 (F1 = \TransNetShotFOne) khoảng $+4.1$ điểm phần trăm trên cùng tập kiểm tra và cùng giao thức đánh giá. Việc để NAS tự khám phá kiến trúc trong không gian 83.9 triệu cấu hình ứng viên (gồm bốn biến thể DDCNN: V2, A, B, C) đem lại lợi thế định lượng rõ rệt so với mạng xương sống cố định được thiết kế thủ công. Lợi thế này được giải thích bởi khả năng NAS thích nghi với phân phối dữ liệu mục tiêu -- video ngắn có shot dày đặc và hiệu ứng kỹ xảo phức tạp -- thay vì giả định một kiến trúc cố định hoạt động tốt cho mọi miền.

\textbf{RQ2 -- Các kỹ thuật ở tầng phân loại có cải thiện F1 mà không cần thay đổi mạng xương sống?}
Câu trả lời là \textbf{có}. Khi giữ nguyên mạng xương sống AutoShot và chỉ tinh chỉnh tầng phân loại bằng tổ hợp Focal Loss, nhãn many-hot, Gaussian smoothing và temperature scaling, AutoShotV2 đạt F1 = \ASVTwoShotFOne{} trên SHOT -- cao hơn AutoShot gốc $+\ASVTwoVsAutoShotShotPP$ điểm phần trăm và cao hơn mô hình cơ sở TransNetV2 $+\ASVTwoVsTransNetShotPP$ điểm phần trăm. Phát hiện này cho thấy ngay cả khi NAS đã tối ưu kiến trúc, vẫn còn dư địa cải thiện đáng kể chỉ thông qua thiết kế lại tầng đầu ra và hậu xử lý xác suất. Mức cải thiện trên đạt được mà không cần huấn luyện mạng xương sống từ đầu, qua đó tách biệt tác động của nhóm kỹ thuật ở tầng phân loại so với phần mạng xương sống.

\textbf{RQ3 -- Mức cải thiện trên SHOT có duy trì khi đánh giá chéo trên BBC và ClipShots?}
Câu trả lời là \textbf{một phần}. Trên BBC, AutoShotV2 đạt F1 = \ASVTwoBBCFOne, vượt TransNetV2 (0.9620) và AutoShot tự chạy lại (0.9554), cho thấy các kỹ thuật hiệu chỉnh xác suất có thể chuyển dịch tốt sang miền video tài liệu trong thiết lập hiện tại. Tuy nhiên, trên ClipShots, F1 của AutoShotV2 đạt \ASVTwoClipFOne, thấp hơn AutoShot theo báo cáo gốc (0.7870) và AutoShot tự chạy lại (0.7753). Hiện tượng \textit{domain-specific overfitting} này được lý giải bởi ba giả thuyết kỹ thuật trong Chương~\ref{ch:ket_qua_thuc_nghiem}: nhãn many-hot mở rộng $\pm 1$ frame chưa phù hợp với phân phối chuyển cảnh của ClipShots, hiệu chỉnh nhiệt độ tối ưu trên tập xác thực của SHOT bị lệch khi áp dụng sang miền khác, và cấu hình Focal Loss ($\gamma=1.0, \alpha=0.5$) thiên về phân phối nhãn của SHOT. Như vậy, RQ3 chỉ ra một giới hạn quan trọng: tối ưu cho một miền cụ thể không tự động chuyển dịch ổn định sang mọi miền khác -- đây là động lực cho Mục~\ref{sec:huong_phat_trien} về xuyên miền calibration.

\section{Đóng góp của khóa luận}

Trên cơ sở ba câu trả lời ở trên, khóa luận đem lại ba nhóm đóng góp chính:

\begin{enumerate}
    \item \textbf{Hệ thống hóa tri thức (nền cho RQ1):} Cung cấp khảo sát có cấu trúc về dòng chảy nghiên cứu SBD, từ phương pháp cổ điển dựa trên đặc trưng thủ công, đến các kiến trúc học sâu CNN/3D ConvNet/Transformer, và gần đây nhất là hướng NAS. Phân tích định lượng trên ba bộ dữ liệu chuẩn (SHOT, BBC, ClipShots) làm rõ khoảng trống nghiên cứu giữa video dài truyền thống và video ngắn hiện đại; khảo sát và đối chiếu thêm với RAI qua số liệu đã công bố. Cơ sở khảo sát này định vị câu trả lời cho RQ1 trong bối cảnh các mạng xương sống học sâu hiện hành.

    \item \textbf{Đóng góp về phương pháp (giải quyết RQ2):} Phân tích và triển khai quy trình AutoShot/AutoShotV2, trong đó bổ sung nhóm kỹ thuật cải tiến ở tầng phân loại và hậu xử lý (Focal Loss, nhãn many-hot, hiệu chỉnh nhiệt độ, làm mượt Gaussian) mà không thay đổi mạng xương sống gốc --- cô lập đóng góp của tầng phân loại và đem lại cải thiện thực chất chỉ thông qua thiết kế tầng đầu ra.

    \item \textbf{Đóng góp về kết quả thực nghiệm (đối với RQ3):} Trên tập SHOT, AutoShotV2 đạt F1 = \ASVTwoShotFOne, cao hơn AutoShot gốc khoảng $+\ASVTwoVsAutoShotShotPP$ điểm phần trăm và cao hơn mô hình cơ sở TransNetV2 khoảng $+\ASVTwoVsTransNetShotPP$ điểm phần trăm; đồng thời cung cấp phân tích so sánh chéo trên BBC và ClipShots làm rõ mức độ tổng quát hóa --- vượt mô hình cơ sở trên BBC nhưng suy giảm trên ClipShots, chỉ ra giới hạn xuyên miền cần khắc phục trong nghiên cứu tiếp theo.

\end{enumerate}

Tóm lại, khóa luận khẳng định rằng việc kết hợp giữa \textit{kiến trúc mạng xương sống linh hoạt} được tối ưu bởi NAS và \textit{quy trình huấn luyện--hậu xử lý chặt chẽ} ở tầng phân loại là chìa khóa để đạt hiệu suất cao trên các nền tảng nội dung video ngắn có nhịp độ nhanh.

\section{Hạn chế của nghiên cứu}
\label{sec:han_che}

Mặc dù đạt được những kết quả tích cực, khóa luận vẫn tồn tại bốn nhóm hạn chế cần được nêu rõ:

\subsection{Chi phí tính toán của giai đoạn NAS}
Quy trình NAS hai giai đoạn của AutoShot (huấn luyện SuperNet + tìm kiếm bằng Bayesian Optimization) cộng với chưng cất tri thức và Weight Grafting có chi phí tổng cộng khoảng 76--91 giờ trên hệ thống đa GPU V100. Đây là một \textit{thách thức triển khai} hơn là lỗ hổng phương pháp luận: trong phạm vi khóa luận, nhóm thực hiện đã tận dụng các điểm lưu trọng số (điểm lưu trọng số) công khai của AutoShot gốc~\cite{zhu2023autoshot} thay vì tự chạy lại quy trình tìm kiếm, cho phép tập trung tài nguyên vào các cải tiến ở tầng phân loại. Tuy giải pháp này hợp lý cho khóa luận, nó kế thừa giả thiết rằng kiến trúc tìm được trên SHOT (theo tác giả gốc) cũng là điểm khởi đầu phù hợp cho miền video ngắn nói chung --- giả thiết này chưa được kiểm chứng độc lập trong khóa luận và là một giới hạn cần lưu ý khi mở rộng sang các phân phối video ngắn khác.

\subsection{Nhạy cảm với chuyển động cực nhanh và hiệu ứng ánh sáng}
Trong một số trường hợp video ngắn có chuyển động máy quay cực nhanh hoặc hiệu ứng đèn flash mạnh, mô hình AutoShotV2 vẫn xảy ra báo động giả (dương tính giả) do nhầm lẫn giữa biến đổi trong nội bộ shot và chuyển cảnh thật. Đặc biệt, các video game/animation với hiệu ứng kỹ xảo dày đặc tạo ra biến thiên thị giác lớn nhưng không tương ứng với chuyển cảnh thật -- một loại nhiễu cấu trúc khó xử lý chỉ bằng tín hiệu thị giác đơn lẻ.

\subsection{Mất cân bằng dữ liệu chưa được giải quyết triệt để}
Mặc dù đã áp dụng Focal Loss và nhãn many-hot, sự chênh lệch giữa số khung hình nền (không phải ranh giới) và khung hình ranh giới vẫn rất lớn (xấp xỉ 1:100 đến 1:200 tùy bộ dữ liệu). Hậu quả định lượng rõ nhất nằm ở chuyển cảnh dần: trên ClipShots, recall theo loại chuyển cảnh của cấu hình triển khai đạt $0.8991$ với chuyển cảnh đột ngột nhưng chỉ $0.8089$ với chuyển cảnh dần (2.302 chuyển cảnh dần trong ground truth) --- tức khoảng một phần năm chuyển cảnh dần bị bỏ sót, cho thấy Focal Loss và many-hot ở cấu hình hiện tại chưa đủ cho lớp mẫu hiếm và khó này. Các hướng khắc phục khả thi gồm lấy mẫu quá mức (oversampling) các cửa sổ chứa chuyển cảnh dần, khai thác mẫu khó (hard example mining) ở các epoch sau, và tự động tinh chỉnh tham số $\gamma, \alpha$ của Focal Loss theo phân phối nhãn của từng tập dữ liệu --- tất cả chưa được khảo sát đầy đủ trong khóa luận.

\subsection{Khoảng cách miền dữ liệu giữa SHOT và các bộ dữ liệu khác}
Như đã chỉ ra ở RQ3, các kỹ thuật trong AutoShotV2 tối ưu cho phân phối của SHOT có biểu hiện suy giảm khoảng $3.1$ điểm phần trăm trên ClipShots. Khoảng cách này có gốc rễ thống kê rõ ràng: SHOT có độ dài shot trung bình $2.59$~giây với hiệu ứng chuyển cảnh dày đặc, trong khi ClipShots trung bình $15.34$~giây/shot với nội dung mở đa dạng hơn 20 danh mục --- một dịch chuyển phân phối (domain shift) tự nhiên khiến ngưỡng và nhiệt độ tối ưu trên SHOT không còn tối ưu trên ClipShots. Khóa luận đã đề xuất ba giả thuyết kỹ thuật giải thích hiện tượng (nhãn many-hot, hiệu chỉnh nhiệt độ tối ưu trên SHOT, cấu hình Focal Loss) nhưng chưa có thử nghiệm phân rã tách biệt từng yếu tố trên ClipShots --- do đó các kết luận về xuyên miền hiện ở mức định hướng phân tích thay vì khẳng định nguyên nhân. Khả năng tổng quát hóa sang dữ liệu sản xuất nội dung tiếng Việt cũng chưa được kiểm chứng và sẽ được thảo luận ở Mục~\ref{sec:huong_phat_trien}.

\section{Hướng phát triển tương lai}
\label{sec:huong_phat_trien}

Trên cơ sở những phát hiện và hạn chế đã phân tích, sáu hướng phát triển khả thi được đề xuất:

\subsection{Tối ưu hóa real-time và triển khai trên thiết bị}
Nghiên cứu các kỹ thuật nén mô hình (cắt tỉa, lượng tử hóa int8/int4, chưng cất tri thức sang mô hình học trò nhỏ) để đưa AutoShotV2 vào các ứng dụng xử lý video trực tiếp với độ trễ thấp. Mục tiêu cụ thể là đạt thông lượng tối thiểu 30 FPS trên các thiết bị di động phổ thông, mở đường cho ứng dụng SBD trên thiết bị cho các nền tảng video ngắn.

\subsection{Phát hiện đa phương thức (Multimodal SBD)}
Kết hợp tín hiệu âm thanh (tín hiệu âm thanh như sự thay đổi đột ngột của track nhạc nền) và thông tin trong miền nén (vector chuyển động của bộ mã hóa-giải mã H.264/H.265) để tăng cường độ bền vững khi tín hiệu thị giác bị nhiễu hoặc mờ nhòe. Đặc biệt, đối với video game, tín hiệu âm thanh thường ổn định hơn hình ảnh và có thể bổ trợ tốt cho việc loại trừ các biến đổi nội bộ shot do kỹ xảo gây ra.

\subsection{Khả năng tổng quát hóa xuyên miền}
Xây dựng cơ chế thích nghi miền (Domain Adaptation) để mô hình hoạt động ổn định trên nhiều nền tảng khác nhau (TikTok, YouTube Shorts, phim điện ảnh truyền thống) mà không cần huấn luyện lại từ đầu. Bám sát ba giả thuyết xuyên miền đã nêu ở Chương~\ref{ch:ket_qua_thuc_nghiem}, các hướng cụ thể gồm: (i) \textit{hiệu chỉnh nhiệt độ tại thời điểm suy diễn (test-time temperature recalibration)} --- tính lại $T^*$ trên một tập xác thực nhỏ của miền đích trước khi suy diễn để khắc phục lệch logit; (ii) \textit{meta-learning cho các tham số Focal Loss} --- học một phân phối khởi tạo $\gamma, \alpha$ thích nghi nhanh khi gặp miền mới, dựa trên thống kê phân phối nhãn; và (iii) \textit{domain-adversarial training} cho tầng phân loại để học biểu diễn đầu ra ít phụ thuộc vào miền nguồn. Ba hướng này cho phép giữ nguyên mạng xương sống đã được tối ưu hóa bằng NAS mà vẫn chuyển dịch được sang miền mới với chi phí thấp.

\subsection{Phân loại fine-grained loại chuyển cảnh}
Mở rộng bài toán từ phân loại nhị phân (boundary/non-boundary) sang phân loại đa lớp với các nhãn cụ thể như Cut, Dissolve, Fade-in, Fade-out, Wipe. Việc này phục vụ cho các hệ thống biên tập video tự động cao cấp như tự động chú thích loại chuyển cảnh khi nhập video thô, hoặc gợi ý hiệu ứng tương tự cho người sáng tạo nội dung.

\subsection{Tự giám sát và mở rộng dữ liệu}
Khai thác các phương pháp học tự giám sát (contrastive learning, masked frame prediction) trên kho dữ liệu video chưa gán nhãn quy mô lớn để mô hình học được các biểu diễn đặc trưng mạnh hơn mà không phụ thuộc nhiều vào nhãn thủ công. Đây là hướng phù hợp với xu thế chung của học sâu trong thị giác máy tính và có thể giảm chi phí gán nhãn cho việc mở rộng dataset.

\subsection{Mở rộng dataset cho video tiếng Việt}
Xây dựng \textit{SHOT-VN} -- một bộ dữ liệu chuyên biệt cho video ngắn tiếng Việt thu thập từ TikTok Việt Nam, YouTube Shorts của các nhà sáng tạo nội dung Việt -- với hệ thống chú thích theo cùng quy chuẩn của SHOT gốc. Bộ dữ liệu này sẽ cho phép đánh giá khả năng tổng quát hóa của các mô hình hiện hành sang văn hóa và phong cách dựng video đặc thù Việt Nam, đồng thời tạo điểm chuẩn cho các nghiên cứu SBD trong nước.

\section{Lời kết}

Bài toán phát hiện ranh giới cảnh quay tưởng chừng là một bước tiền xử lý đơn giản, nhưng khi đặt trong bối cảnh video ngắn hiện đại với nhịp cắt cảnh dày đặc, hiệu ứng kỹ xảo phức tạp và sự đa dạng cao của nội dung, nó lại trở thành một bài toán có nhiều thách thức kỹ thuật và ý nghĩa thực tiễn. Khóa luận đã đóng góp một bước tiến nhỏ trên hành trình giải quyết bài toán này: từ việc hệ thống hóa khảo sát, đến đề xuất AutoShotV2 với F1 = \ASVTwoShotFOne{} trên SHOT, và quan trọng không kém -- chỉ ra rõ những giới hạn cần khắc phục về khả năng tổng quát hóa đa miền.

Trong thời đại video ngắn tiếp tục bùng nổ trên các nền tảng số, các hệ thống phân tích video tự động chính xác và hiệu quả sẽ là nền tảng cho nhiều ứng dụng giá trị -- từ kiểm duyệt nội dung, gợi ý cá nhân hóa, đến công cụ biên tập sáng tạo. Khóa luận này chỉ ra rõ giới hạn của khả năng tổng quát hóa đa miền trong các phương pháp SBD hiện hành và mở đường cho hai hướng nghiên cứu cụ thể: (i) cơ chế hiệu chỉnh xác suất thích nghi theo miền đích thay vì cố định theo miền nguồn, và (ii) xây dựng bộ dữ liệu SHOT-VN cho video ngắn tiếng Việt làm bệ phóng cho các nghiên cứu SBD trong nước. Nhóm thực hiện hy vọng các phân tích và quy trình AutoShotV2 sẽ là điểm tham chiếu hữu ích cho các công trình tiếp theo của cộng đồng học thuật trong và ngoài nước.
